\chapter{Wound Care (Wound Management)}

\section{Definition}
Wound care involves the assessment, cleaning, and management of skin wounds to prevent infection and promote healing. Wounds range from minor cuts and abrasions to complex surgical wounds and chronic ulcers.

\section{Pathophysiology}
Wound healing occurs in four phases: hemostasis (clot formation), inflammation (immune response), proliferation (new tissue formation), and remodeling (scar maturation). Proper wound care optimizes these phases by keeping the wound clean, moist, and protected. Factors that impair healing include infection, ischemia, foreign bodies, malnutrition, and chronic diseases.

\section{Etiology}
\begin{itemize}
\item Acute wounds: cuts, lacerations, abrasions, punctures, surgical incisions
\item Chronic wounds: pressure sores, diabetic foot ulcers, venous stasis ulcers, arterial ulcers
\item Thermal wounds: burns, frostbite
\item Traumatic wounds: avulsions, degloving injuries, crush injuries
\item Infected wounds: abscesses, cellulitis with drainage
\end{itemize}

\section{Indications for Evaluation}
\begin{itemize}
\item Any break in the skin requiring cleaning and protection
\item Wounds with active bleeding that requires control
\item Wounds with dirt, debris, or foreign material
\item Signs of infection: redness, warmth, swelling, pus, red streaks
\item Chronic or non-healing wounds
\item Surgical wounds with drains or staples
\end{itemize}

\section{Related Symptoms}
\begin{itemize}
\item Scar formation
\item Cellulitis (wound infection)
\item Tetanus (risk with contaminated wounds)
\item Abscess formation
\item Impetigo (superficial infection)
\end{itemize}

\section{Self-Care/Home Management}
\begin{itemize}
\item Clean the wound gently with mild soap and water or saline
\item Remove debris with tweezers cleaned with alcohol
\item Apply antibiotic ointment and cover with sterile bandage
\item Change dressing daily or when wet/soiled
\item Keep the wound moist but not wet
\item Do not pick at scabs
\item Monitor for signs of infection
\end{itemize}

\section{Diagnostic Approach}
\begin{itemize}
\item Clinical assessment of wound size, depth, and characteristics
\item Wound culture if infection suspected
\item X-ray to rule out foreign body or underlying fracture
\item Ankle-brachial index for leg wounds to assess arterial supply
\end{itemize}

\section{Treatment/Management Overview}
\begin{itemize}
\item Cleaning and debridement of devitalized tissue
\item Appropriate wound dressing selection (hydrocolloid, foam, alginate)
\item Negative pressure wound therapy for complex wounds
\item Antibiotics for infected wounds
\item Tetanus prophylaxis (tetanus toxoid or immune globulin)
\item Surgical closure for appropriate wounds
\item Management of underlying cause (offloading for pressure sores, compression for venous ulcers)
\end{itemize}
